% 2010 年湖南卷文科第 8 题选项图：四个图象（A）—（D）。
% 来源 PDF：普通高考/2010/2010湖南文.pdf，SHA-256 38611c33c9da2139ad6cd4117513ee144b09a6860180eb976c446b4858ada122。
% 原页：第 1 页（A3 三栏）第 1 栏下部，书页码 603。
% 重绘依据：照原页四个小图——两曲线为 y=ax^2+bx（过原点）与 y=log_{|b/a|}x（只画 x>0）；
%   (A) 抛物线开口向上、另一根在 0 与 1 之间，对数曲线单增；
%   (B) 抛物线开口向下、另一根在 0 与 1 之间，对数曲线单增；
%   (C) 抛物线开口向下、另一根小于 -1，对数曲线单减；
%   (D) 抛物线开口向上、另一根在 -1 与 0 之间，对数曲线单减。
%   各曲线为示意曲线，仅保持原页的开口方向、根与刻度的相对位置及对数曲线单调性。
\begin{tikzpicture}[scale=0.78, semithick]

  % ============ (A) ============
  \begin{scope}
    \begin{scope}
      \clip (-1.6,-1.2) rectangle (2.55,2.4);
      \draw plot[domain=-1.15:1.45, samples=90] (\x,{2.0*\x*(\x-0.95)});
      \draw plot[domain=0.55:2.6, samples=90] (\x,{ln(\x)/ln(2)});
    \end{scope}
    \draw[->] (-1.65,0) -- (2.65,0) node[right] {$x$};
    \draw[->] (0,-1.2) -- (0,2.45) node[above] {$y$};
    \draw (-1,0.09) -- (-1,-0.09);\node[below] at (-1,-0.18) {$-1$};
    \draw (1,0.09) -- (1,-0.09);\node[below] at (1,-0.18) {$1$};
    \node[below left] at (0,-0.10) {$O$};
    \node at (0.5,-1.75) {（A）};
  \end{scope}

  % ============ (B) ============
  \begin{scope}[shift={(10.2,0)}]
    \begin{scope}
      \clip (-1.6,-1.2) rectangle (2.55,2.4);
      \draw plot[domain=-1.1:1.35, samples=90] (\x,{-2.0*\x*(\x-0.9)});
      \draw plot[domain=0.55:2.6, samples=90] (\x,{ln(\x)/ln(2)});
    \end{scope}
    \draw[->] (-1.65,0) -- (2.65,0) node[right] {$x$};
    \draw[->] (0,-1.2) -- (0,2.45) node[above] {$y$};
    \draw (-1,0.09) -- (-1,-0.09);\node[below] at (-1,-0.18) {$-1$};
    \draw (1,0.09) -- (1,-0.09);\node[below] at (1,-0.18) {$1$};
    \node[above left] at (0,0.02) {$O$};
    \node at (0.5,-1.75) {（B）};
  \end{scope}

  % ============ (C) ============
  \begin{scope}[shift={(0,-6.9)}]
    \begin{scope}
      \clip (-1.6,-1.2) rectangle (2.55,2.4);
      \draw plot[domain=-1.8:0.55, samples=90] (\x,{-2.0*\x*(\x+1.3)});
      \draw plot[domain=0.4:2.6, samples=90] (\x,{-ln(\x)/ln(2)});
    \end{scope}
    \draw[->] (-1.65,0) -- (2.65,0) node[right] {$x$};
    \draw[->] (0,-1.2) -- (0,2.45) node[above] {$y$};
    \draw (-1,0.09) -- (-1,-0.09);\node[below] at (-1,-0.18) {$-1$};
    \draw (1,0.09) -- (1,-0.09);\node[below] at (1,-0.18) {$1$};
    \node[below left] at (0,-0.10) {$O$};
    \node at (0.5,-1.75) {（C）};
  \end{scope}

  % ============ (D) ============
  \begin{scope}[shift={(10.2,-6.9)}]
    \begin{scope}
      \clip (-1.6,-1.2) rectangle (2.55,2.4);
      \draw plot[domain=-1.5:1.25, samples=90] (\x,{2.0*\x*(\x+0.9)});
      \draw plot[domain=0.4:2.6, samples=90] (\x,{-ln(\x)/ln(2)});
    \end{scope}
    \draw[->] (-1.65,0) -- (2.65,0) node[right] {$x$};
    \draw[->] (0,-1.2) -- (0,2.45) node[above] {$y$};
    \draw (-1,0.09) -- (-1,-0.09);\node[below] at (-1,-0.18) {$-1$};
    \draw (1,0.09) -- (1,-0.09);\node[below] at (1,-0.18) {$1$};
    \node[below right] at (0,-0.06) {$O$};
    \node at (0.5,-1.75) {（D）};
  \end{scope}

\end{tikzpicture}
